\documentclass[11pt]{article}
\usepackage[margin=0.25in]{geometry}
\usepackage{array}
\usepackage{booktabs}
\usepackage{caption}
\usepackage{graphicx}
\usepackage{hyperref}
\usepackage{longtable}
\usepackage{lmodern}
\usepackage{makecell}
\usepackage[expansion=false]{microtype}
\usepackage{ragged2e}
\usepackage{tabularx}
\usepackage{xurl}
\pagestyle{empty}
\begin{document}
\begingroup
\small
\setlength{\tabcolsep}{4pt}
\begin{center}
\begin{tabularx}{0.98\linewidth}{@{}>{\RaggedRight\arraybackslash}p{0.10\linewidth}>{\RaggedRight\arraybackslash}p{0.18\linewidth}>{\RaggedRight\arraybackslash}X>{\RaggedRight\arraybackslash}X@{}}
\toprule
Method & Model & Primary role & Main trade-off \\
\midrule
A & Gaussian VaR & Thin-tailed rolling benchmark used to test whether a Gaussian assumption is adequate. & Transparent baseline, but structurally weak under heavy tails and volatility clustering. \\
B & Student-t VaR & Parametric heavy-tail benchmark using a rolling return window. & Transparent and fast, but only weakly conditional on current market state. \\
C & Jump-diffusion Monte Carlo & Scenario engine that combines continuous diffusion with discrete jump risk. & More realistic under stress, but sensitive to jump-calibration choices. \\
D & LSTM Gaussian head & Feature-conditioned one-step forecast of mean and volatility. & Adaptive to regime shifts, but harder to validate and explain. \\
E & VAE latent sampling & Generative model for synthetic return scenarios and latent tail structure. & Useful for scenario enrichment, but weaker as a standalone production VaR engine. \\
\bottomrule
\end{tabularx}
\end{center}
\endgroup
\end{document}
